\chapter{Совместный квадрупольно-тетраэдрический вакуум и остаточная группа
\texorpdfstring{$Z_3$}{Z3}}

Предыдущие главы отдельно вывели поле квадрупольного порядка
$Q\in\operatorname{Sym}_0^2(V_1)$, тетраэдрическое поле
$T\in\operatorname{Sym}_0^3(V_1)$ и одну семейную связность $B$. Однако
простая сумма двух потенциалов не задаёт их взаимную ориентацию: она имеет
пять орбитальных нулей вместо трёх нулей общего $SO(3)$. Здесь строится
смешанная кривизна, которая обязана сцепить обе фазы без нового
подбираемого коэффициента.

\section{Антикруговой журнал и литературная граница}

В общей теории тензорных фаз совместное присутствие параметров порядка
второго и третьего рангов допустимо и необходимо для описания смешанных
нематико-тетраэдрических фаз; систематическая классификация дана в
\path{arXiv:cond-mat/0205171}. Нарушение $SO(3)\to A_4$ полем спина три
использовано в \path{arXiv:0910.4392}, а конечные вихри вакуумного
многообразия $SO(3)/A_4$ построены в \path{arXiv:2607.12366}.

Эти источники подтверждают представительный тип и физическую осмысленность
совместного порядка, но не дают следующую проектную формулу и не фиксируют
её коэффициент. Поэтому смешанная кривизна ниже является новым внутренним
результатом проекта, а не цитатой из литературы.

\section{Проективное чтение поля \texorpdfstring{$Q$}{Q}}

Пусть $\Delta=a_*-b_*$ --- уже выведенная щель упорядоченного спектра
$R_*=\operatorname{diag}(a_*,b_*,b_*)$. На упорядоченной орбите
\begin{equation}
 Q=\Delta\left(P-\frac13I_3\right),
 \qquad
 P_Q=\frac13I_3+\frac{Q}{\Delta}=P=nn^{\mathsf T}.
 \label{eq:v6-coupled-projective-readout}
\end{equation}
Таким образом, $P_Q$ читает не ориентированный вектор $n$, а его
проективную ось. Это существенно: прежнее вакуумное многообразие поля
$Q$ было $\mathbb{RP}^2$, поэтому смешанный член не должен выбирать знак
директора вручную.

Свернём тетраэдрический тензор с этим проектором:
\begin{equation}
 w_i=T_{ijk}(P_Q)_{jk}.
 \label{eq:v6-coupled-composed-arrow}
\end{equation}
Поскольку $T_{ijj}=0$, та же формула имеет вид
$w_i=T_{ijk}Q_{jk}/\Delta$. Она является композицией уже существующих
носителей $V_2\to V_1$, а не новым независимым полем.

\section{Смешанная матричная кривизна}

Для канонического тетраэдра с вершинами $n_a$ имеем
\begin{equation}
 T_*=\sum_{a=1}^{4}n_a^{\otimes3}.
\end{equation}
Если $P=n_1n_1^{\mathsf T}$, то
\begin{equation}
 T_*:P=\frac89n_1.
 \label{eq:v6-tetrahedral-projector-contraction}
\end{equation}
Число $8/9$ не является свободной константой: оно следует из
$n_a\cdot n_b=-1/3$ при $a\ne b$. Поэтому определим
\begin{equation}
 \mathcal F_{QT}
 =ww^{\mathsf T}-\frac{64}{81}P_Q,
 \qquad
 V_{QT}=\frac13\Tr\mathcal F_{QT}^2.
 \label{eq:v6-coupled-mixed-curvature}
\end{equation}
Это один положительный матричный квадрат на том же триплетном углу.
Нового относительного веса или нового массового масштаба не введено.

На выровненном вакууме остаток
$\|\mathcal F_{QT}\|=5.9\cdot10^{-16}$. Норма смешанного блока полного
гессиана равна $1.761249\ldots$, поэтому член действительно связывает
$Q$ и $T$, а не распадается на переименование двух независимых
потенциалов.

\begin{proposition}[Точное множество выбранных осей]
При фиксированном каноническом $T_*$ и проекторе $P=nn^{\mathsf T}$
условие $\mathcal F_{QT}=0$ выполняется тогда и только тогда, когда
\begin{equation}
 |n_1|=|n_2|=|n_3|=\frac1{\sqrt3}.
\end{equation}
Следовательно, существуют восемь ориентированных решений, то есть ровно
четыре проективные тетраэдрические оси.
\end{proposition}

Действительно, в координатах канонического тетраэдра
$w\propto(n_2n_3,n_1n_3,n_1n_2)$. Равенство двух матриц ранга один
сначала требует их общего образа, а затем равенства норм; при ненулевом
$n$ это даёт равенство модулей трёх координат. Прямое перечисление восьми
знаковых троек даёт максимальный остаток $1.0\cdot10^{-15}$.

\section{Полный потенциал и устранение изотропной ветви}

Совместный однородный потенциал имеет вид
\begin{equation}
 V(Q,T)=
 \bigl(\mathcal F_{\beta_c}(I_3/3+Q)-\mathcal F_{\beta_c}(R_*)\bigr)
 +\frac13\Tr\mu_T(T)^2
 +\frac13\Tr\mathcal F_{QT}(Q,T)^2.
 \label{eq:v6-coupled-total-potential}
\end{equation}
Все три слагаемых уже имеют фиксированную нормировку. На выровненном
упорядоченном вакууме их значения равны нулю с точностью
$9.5\cdot10^{-16}$.

При $Q=0$ и $T=T_*$ первые два слагаемых также находятся на прежних
минимумах сосуществования, но смешанный квадрат даёт
\begin{equation}
 V_{QT}(0,T_*)=\frac{4096}{59049}=0.06936611966\ldots>0.
\end{equation}
Тем самым вырожденная изотропная ветвь, сохранявшаяся у одного
канонического потенциала $Q$, больше не является нулём полного действия.

\section{Полный смешанный гессиан}

До сцепления двенадцатимерный гессиан имеет пять численных нулей: три
общих вращения и два относительных поворота оси $Q$ внутри фиксированного
тетраэдра. После добавления~\eqref{eq:v6-coupled-mixed-curvature} остаются
ровно три нуля общей $SO(3)$-орбиты.

После факторизации по этим трём направлениям физический гессиан имеет
девять собственных значений
\begin{equation}
 \begin{split}
 \operatorname{Spec}H_{\mathrm{phys}}=\{&
 1.21473323^{(2)},\ 1.47985517,\
 2.37881858^{(2)},\ 3.39901548,\\
 &6.55303887,\ 19.70072077^{(2)}\}.
 \end{split}
 \label{eq:v6-coupled-physical-hessian-spectrum}
\end{equation}
Все девять мод положительны. Малые значения порядка $10^{-7}$ в полном
гессиане лежат только на трёх точных калибровочных касательных и являются
погрешностью конечных разностей; норма действия гессиана на
нормированную калибровочную орбиту равна $9.1\cdot10^{-7}$.

\begin{theorem}[Локальная устойчивость совместного вакуума]
Потенциал~\eqref{eq:v6-coupled-total-potential} имеет локально устойчивый
совместный вакуум $Q_*+T_*$: после факторизации общей калибровочной
орбиты его гессиан строго положителен. Смешанная кривизна снимает ровно
две относительные ориентационные моды и не создаёт отрицательных мод.
\end{theorem}

\section{Семейная связность и остаточная группа}

Кинетическая массовая матрица одной семейной связности получает вклады
от обоих конденсатов. Их спектры равны
\begin{align}
 \operatorname{Spec}M_Q&=\{0,1.50771578,1.50771578\},\\
 \operatorname{Spec}M_T&=\{14.22222222^{(3)}\},\\
 \operatorname{Spec}(M_Q+M_T)&=
 \{14.22222222,15.72993800^{(2)}\}.
\end{align}
Следовательно, непрерывного стабилизатора нет и все три компоненты
семейной связности массивны.

Из двенадцати собственных вращений тетраэдра ровно три сохраняют
выбранную проективную ось. Они образуют группу $Z_3$, поэтому
\begin{equation}
 SO(3)\longrightarrow A_4\longrightarrow Z_3,
 \qquad
 \mathcal M_{\mathrm{vac}}=SO(3)/Z_3.
\end{equation}

\section{Вердикт и открытая граница}

\begin{protocol}
\begin{enumerate}
 \item смешанная кривизна из прежних носителей: \textbf{получена};
 \item новый относительный коэффициент: \textbf{нет};
 \item четыре проективные тетраэдрические оси: \textbf{выбраны точно};
 \item лишние относительные нулевые моды: \textbf{сняты};
 \item физический гессиан: \textbf{девять положительных мод};
 \item отрицательные физические моды: \textbf{нет};
 \item все три компоненты семейной связности: \textbf{массивны};
 \item остаточная калибровочная группа: \textbf{$Z_3$};
 \item рождение наблюдаемой материи: \textbf{не выведено};
 \item следующий гейт: \textbf{пространственный профиль совместного
 дефекта $Q+T+B$}.
\end{enumerate}
\end{protocol}

Главное изменение статуса модели состоит в том, что её бозонный вакуум
теперь не является набором несвязанных предположений. Поле $Q$,
тетраэдрический порядок $T$ и одна связность $B$ образуют единый
устойчивый узел с остаточной $Z_3$. Но этот результат ещё не доказывает,
что конечные дефекты имеют нужные заряды, фермионные нулевые моды и
спектр Стандартной модели. Следующий тест обязан построить конечную
пространственную конфигурацию, а не переходить к феноменологической
подгонке масс.

Машинный сертификат сохранён в
\path{s2t_v6_bosonic_defect_q_tetrahedral_coupled_vacuum_gate_results.json}.